\documentclass[10pt,a4paper]{amsart}
\usepackage[margin=19mm]{geometry}
\usepackage{amsmath,amssymb,mathtools}
\usepackage[scr=rsfs,cal=euler]{mathalfa}
\usepackage{xfrac}
\usepackage[unicode,hidelinks,bookmarks=false]{hyperref}
\newcommand{\ie}{i.e.\ }
\newcommand{\cf}{cf.\ }
\newcommand{\resp}{respectively\ }
\newcommand{\Subsection}{\subsection}
\providecommand{\nobreakdash}{\nobreak\hbox{-}}
\setlength{\emergencystretch}{2em}
\multlinegap0pt
\usepackage{bm}
\usepackage{bbm}
\usepackage{dsfont}
\usepackage{xspace}
\usepackage{cancel}
\usepackage{mathtools}
\usepackage{amsthm}
\makeatletter
\@ifundefined{thm}{\newtheorem{thm}{Thm}}{}
\@ifundefined{defn}{\newtheorem{defn}[thm]{Definition}}{}
\@ifundefined{prop}{\newtheorem{prop}[thm]{Proposition}}{}
\makeatother
\newcommand{\e}{\mathrm{e}}
\renewcommand{\i}{\mathrm{i}}
\title[The $k$th Order Preserving Sets and Isoperimetric Type Inequ]{The $k$th Order Preserving Sets and Isoperimetric Type Inequalities for Planar Ovals}
\author{Maksymilian Filip Safarewicz, Micha{\l} Zwierzy\'nski}
\date{Source: arXiv:2505.08017v3 (CC BY 4.0)}
\begin{document}
\maketitle

\noindent\emph{Source and licence.} The $k$th Order Preserving Sets and Isoperimetric Type Inequalities for Planar Ovals. Version arXiv:2505.08017v3 (CC BY 4.0), distributed under the Creative Commons Attribution 4.0 International licence, as recorded in the arXiv licence metadata for this exact version and shown on the abstract page https://arxiv.org/abs/2505.08017v3. Full rights notice: licenses/arxiv-2505.08017-CC-BY-4.0.txt.\par\smallskip

\noindent\textit{Editorial scope.} Counted unit: the Prop whose source label is \texttt{Preservingsuppth} in the article (source0.tex lines 349--356, source label \texttt{Preservingsuppth}), delivered together with the article's own complete original proof of it (source0.tex lines 357--382, one \texttt{proof} environment of 26 lines). No mathematics is written, repaired or re-derived: statement and proof are the article's own text line for line. Required context retained uncounted: eqHedghehogForm (lines 169--172); param (lines 192--194); k-tuple pair (lines 257--273). Every definition, display and label of those lines is retained unchanged. Omitted from this package and not redistributed: 1--168, 173--191, 195--256, 274--348, 383--1109. Nothing inside a retained excerpt is omitted or altered: every retained body is delivered exactly as the article's lines are written, so no omission receipt of any kind accompanies this package. Citations: the retained excerpts carry no citation command at all, so no bibliography entry of the article is transcribed and no reference is redistributed. Reference adaptation: none was needed and none was made; every reference inside the delivered unit resolves inside it, so no unresolved reference or citation is produced. Printed slips kept literal: no printed word, symbol, hypothesis or number of the counted statement or of its proof was altered. Layout and class: the article's own class is \texttt{amsart} with packages including graphicx, amsaddr, amssymb, amsmath, amsthm, subcaption, enumerate, xcolor, MnSymbol, bbm, calligra, hyperref; the wrapper is the lane's pinned \texttt{amsart} editorial wrapper at 10pt on A4, which restates the theorem-like environments the retained text needs and adds the article's own macro definitions that the retained text needs (\texttt{\textbackslash e}, \texttt{\textbackslash i}), with extra wrapper packages (bm, bbm, dsfont, xspace, cancel, mathtools). The delivered numbering is the unit's own. Assets: no retained span contains a figure environment, a graphics-inclusion command, a PDF-page inclusion, a table or a TikZ picture; no image file is copied into this package, the package declares no asset dependency and no image file is redistributed with it. Licence: version arXiv:2505.08017v3 of this article is distributed under the Creative Commons Attribution 4.0 International licence, as recorded in the arXiv licence metadata for this exact version and shown on the abstract page https://arxiv.org/abs/2505.08017v3. A full rights notice is delivered at licenses/arxiv-2505.08017-CC-BY-4.0.txt. Characters: every retained excerpt is pure ASCII, so the delivered engine, XeLaTeX with its default Latin Modern text font, reports no missing character.
\begin{align}
\label{eqHedghehogForm}\mathcal{H}(s) &= h(s)\, u(s) + h'(s)\, u'(s)\\ 
\nonumber&=\Big( h(s) \cos(s) - h'(s) \sin(s),  h(s) \sin(s) + h'(s) \cos(s) \Big),
\end{align}

\begin{align}
    \label{param}\mathcal{H}^\perp(s) = -\Big( h(s) \sin(s) + h'(s) \cos(s),  -h(s) \cos(s) + h'(s) \sin(s) \Big).
\end{align}

\begin{defn}\label{k-tuple pair}
Let $\mathcal{H}$ be a hedgehog and let $k>2$ be an integer. We call points
$p_1,\ldots,p_k\in\mathcal{H}$ an \textit{isogonal family of points} if, for each
$i=1,\ldots,k$ (with the convention $p_{k+1}=p_1$), the support lines to
$\mathcal{H}$ at $p_i$ and $p_{i+1}$ meet at the same angle, that is, $\frac{2\pi}{k}$.
Equivalently, the unit normals (compatible with the orientation of $\mathcal{H}$) at $p_1,\ldots,p_k$ are spaced by angles~$\frac{2\pi}{k}$.

Let $\mathcal{H}^{\perp}$ be the image of $\mathcal{H}$ under a counterclockwise
rotation by $\frac{\pi}{2}$ about the origin. Consider the parameterization
of $\mathcal{H}$ given by equation~\eqref{eqHedghehogForm} and the induced
parameterization of $\mathcal{H}^{\perp}$ as in~\eqref{param}. For $s\in[0,2\pi]$,
write $p=\mathcal{H}(s)$ and define $p^{\perp}=\mathcal{H}^{\perp}(s)$. We denote by
$p_1^{\perp},\ldots,p_k^{\perp}\in\mathcal{H}^{\perp}$ the images of an isogonal family
$p_1,\ldots,p_k\in\mathcal{H}$; in particular, for every $i$, the support line to
$\mathcal{H}^{\perp}$ at $p_i^{\perp}$ is perpendicular to the support line to
$\mathcal{H}$ at $p_i$.
\end{defn}

\begin{prop}\label{Preservingsuppth}
Let $\mathcal{H}$ be a hedgehog with the support function $h: [0,2\pi] \to \mathbb{R}$, where $h(s)=a_0+\sum_{n=1}^{\infty}a_n\cos (ns)+b_n\sin (ns)$ is its Fourier series representation, and let the average width of $\mathcal{H}$ be $\overline{w}$. Then, its $k$th Order Preserving Set is a~hedgehog with a support function given by the formula:
$$ h_k(s) = \frac{1}{k}\sum_{j=0}^{k-1}h\left(s+\frac{2\pi j}{k}\right) - \frac{1}{2}\overline{w},$$
and its Fourier series representation is as follows:
\begin{align}
    \label{eqFourierkOPS}h_k(s) = \sum_{k \mid n,\ n > 1}^\infty a_n\cos(ns) + b_n\sin(ns).
\end{align}
\end{prop}

\begin{proof}
 Note that the points $\mathcal{H}(s),\ldots,\mathcal{H}\left(s+\frac{2\pi(k-1)}{k}\right)$, for a given $s\in[0,2\pi]$, constitute an isogonal family of $\mathcal{H}$. Then, the parameterization of the $k$th Order Preserving Set is:
    $$\mathcal{P}_k(s) =\frac{1}{k}\sum_{j=1}^{k}\Bigg(\cos\left(\frac{2\pi j}{k}\right)\cdot \mathcal{H}\left(s+\frac{2\pi j}{k}\right)-\sin\left(\frac{2\pi j}{k}\right)\cdot \mathcal{H}^\perp\left(s+\frac{2\pi j}{k}\right)\Bigg) -\frac{1}{2}\overline{w}\cdot u(s).$$\\Naturally, the point $\mathcal{H}^\perp(s+\frac{2\pi j}{k})$ is understood as in Definition \ref{k-tuple pair}. Using the parameterization \eqref{param} of $\mathcal{H}^\perp$ we obtain the following equalities:
    {\footnotesize
    \begin{align*}    
    \mathcal{P}_k&(s) = \frac{1}{k}\sum_{j=1}^{k}\Bigg(\cos\Big(\frac{2\pi j}{k}\Big)\cdot \Bigg( h\Big(s+\frac{2\pi j}{k}\Big)\cos\Big(s+\frac{2\pi j}{k}\Big) - h'\Big(s+\frac{2\pi j}{k}\Big) \sin\Big(s+\frac{2\pi j}{k}\Big),\\&  h\Big(s+\frac{2\pi j}{k}\Big) \sin\Big(s+\frac{2\pi j}{k}\Big) + h'\Big(s+\frac{2\pi j}{k}\Big) \cos\Big(s+\frac{2\pi j}{k}\Big) \Bigg)+\\&+\frac{1}{k}\sum_{j=1}^{k}\Bigg(\sin\Big(\frac{2\pi j}{k}\Big)\cdot \Bigg( h\Big(s+\frac{2\pi j}{k}\Big) \sin\Big(s+\frac{2\pi j}{k}\Big) + h'\Big(s+\frac{2\pi j}{k}\Big) \cos\Big(s+\frac{2\pi j}{k}\Big),\\&   -h\Big(s+\frac{2\pi j}{k}\Big) \cos\Big(s+\frac{2\pi j}{k}\Big) + h'\Big(s+\frac{2\pi j}{k}\Big) \sin\Big(s+\frac{2\pi j}{k}\Big) \Bigg)\Bigg)-\frac{1}{2}\overline{w}\cdot \big(\cos(s),\sin(s)\big)\\&=\frac{1}{k}\sum_{j=1}^{k}\Bigg( \Bigg( h\Big(s+\frac{2\pi j}{k}\Big)\cos(s) - h'\Big(s+\frac{2\pi j}{k}\Big) \sin(s),\\&  h\Big(s+\frac{2\pi j}{k}\Big) \sin(s) + h'\Big(s+\frac{2\pi j}{k}\Big) \cos(s) \Bigg)-\frac{1}{2}\overline{w}\cdot \big(\cos(s),\sin(s)\big)\\&= \Big(\cos(s),\sin(s)\Big)\cdot \Bigg(\frac{1}{k}\sum_{j=1}^{k} h\Big(s+\frac{2\pi j}{k}\Big)-\frac{1}{2}\overline{w}\Bigg) + \Big(-\sin(s),\cos(s)\Big)\cdot \Bigg(\frac{1}{k}\sum_{j=1}^{k} h'\Big(s+\frac{2\pi j}{k}\Big)\Bigg).
    \end{align*}}


    
    Now, setting $h_k(s) = \frac{1}{k}\sum_{j=0}^{k-1}h\left(s+\frac{2\pi j}{k}\right) - \frac{1}{2}\overline{w}$, we obtain $P_k(s)$ in the form of equation \eqref{eqHedghehogForm}. Therefore, we end the first part of the proof. To derive equation \eqref{eqFourierkOPS}, one can simplify the aforementioned form of $h_k$ by leveraging the properties of Fourier and geometric series. Let us consider the Fourier expansion of $h$:
    $$h(s) = \sum_{n=-\infty}^\infty c_n{\e}^{\i ns}.$$
    Hence:
     $$h\left(s+\frac{2\pi j}{k}\right) = \sum_{n=-\infty}^\infty c_n{\e}^{\i n\left(s+\frac{2\pi j}{k}\right)}.$$
     This series is convergent. Thus, the following is well-defined:
     $$\sum_{j=0}^{k-1}h\left(s+\frac{2\pi j}{k}\right) = \sum_{j=0}^{k-1}\sum_{n=-\infty}^\infty c_n{\e}^{\i n\left(s+\frac{2\pi j}{k}\right)}=\sum_{n=-\infty}^\infty \sum_{j=0}^{k-1} c_n{\e}^{\i n\left(s+\frac{2\pi j}{k}\right)}.$$ Notice that $\sum_{n=-\infty}^\infty \sum_{j=0}^{k-1} c_n{\e}^{\i n\left(s+\frac{2\pi j}{k}\right)} = \sum_{n=-\infty}^\infty c_n{\e}^{\i ns}\sum_{j=0}^{k-1} {\e}^{\i\frac{2\pi jn}{k}}$. We shall now simplify the geometric series:
     $$\sum_{j=0}^{k-1} {\e}^{\i\frac{2\pi jn}{k}} = \sum_{j=0}^{k-1} \left({\e}^{\i\frac{2\pi n}{k}}\right)^j = \begin{cases} 
k & \text{if } k\mid n, \\ 
0 & \text{if } k\nmid n. 
\end{cases}$$ Therefore, we arrive at:
$$\sum_{j=0}^{k-1}h\left(s+\frac{2\pi j}{k}\right) = k\sum_{k\mid n} c_n{\e}^{\i ns}.$$ Translating the result into trigonometric series one gets that:
\begin{align*}
\frac{1}{k}\sum_{j=0}^{k-1}h\left(s+\frac{2\pi j}{k}\right) - \frac{1}{2}\overline{w} &= \sum_{k \mid n} a_n\cos(ns) + b_n\sin(ns) - a_0\\ &= \sum_{k \mid n,\ n > 1}a_n\cos(ns) + b_n\sin(ns),
\end{align*}
which concludes the proof.
    \end{proof}

\end{document}
